\subsection{Berechnung der Ausgangsspannung mit mehreren Eingangsspannungen}
% Alt: Aufgabe 14
\Aufgabe{
    \label{Aufg14}
    Gegeben ist die abgebildete Schaltung mit einem idealen Operationsverstärker. Die Betriebsspannung
    beträgt $\pm 15\,\mathrm{V}$.
    \begin{itemize}
        \item[a)] Die Widerstände sind alle gleich, es gilt $R_1 = R_2 = R_3 = R_\mathrm{F} = \mathrm{1\,k\Omega}$.
              Die Spannung $U_\mathrm{E1} = \mathrm{1\,V}$, $U_\mathrm{E2} = \mathrm{2\,V}$, $U_\mathrm{E3} = \mathrm{3\,V}$.
              Wie hoch ist die Spannung am Ausgang $U_\mathrm{A}$?
        \item[b)] Die Eingangsspannungen bleiben gleich groß, nur die Widerstände betragen nun
              $R_\mathrm{F} = \mathrm{10\,k\Omega}, R_1=\mathrm{2\,k\Omega}, R_2=\mathrm{5\,k\Omega}, R_3=\mathrm{10\,k\Omega}$.
              Wie hoch ist die Spannung am Ausgang U$_A$?
        \item[c)] Die Eingangsspannungen bleiben gleich groß, nur die Widerstände betragen nun
              $R_\mathrm{F} = \mathrm{10\,k\Omega}, R_1 = \mathrm{10\,k\Omega}, R_2 = \mathrm{5\,k\Omega}, R_3 = \mathrm{2\,k\Omega}$.
              Wie hoch ist die Spannung am Ausgang $U_\mathrm{A}$?
    \end{itemize}
    \begin{figure}[H]
        \centering
        \input{Tikz/src/FigOPVmehrereEingaenge}\alttext{Das Bild zeigt eine Schaltung mit einem Operationsverstärker. Auf der linken Seite sind drei Eingangsspannungen \(U_{\mathrm{E1}}\), \(U_{\mathrm{E2}}\) und \(U_{\mathrm{E3}}\), die jeweils über Widerstände \(R_1\), \(R_2\) und \(R_3\) zum invertierenden Eingang des Operationsverstärkers führen. Der nicht-invertierende Eingang ist mit Masse verbunden. Vom Ausgang des Operationsverstärkers führt ein Rückkoppelwiderstand \(R_\mathrm{F}\) zurück zum invertierenden Eingang. Die Ausgangsspannung ist mit \(U_\mathrm{A}\) bezeichnet. Alle Spannungen sind durch Pfeile von oben nach unten oder von links nach rechts gekennzeichnet.}
        \label{fig:FigAusgangsspannungmehrereEingaenge}
    \end{figure}
}

\Loesung{

    Gegeben ist eine Summierschaltung mit einem idealen Operationsverstärker und einer Betriebsspannung von $\pm 15\,\mathrm{V}$.

    Die allgemeine Formel für die Ausgangsspannung lautet:
    \[
        U_\mathrm{A} = -\left( U_\mathrm{E1} \cdot \frac{R_\mathrm{F}}{R_1} + U_\mathrm{E2} \cdot \frac{R_\mathrm{F}}{R_2} + U_\mathrm{E3} \cdot \frac{R_\mathrm{F}}{R_3} \right)
    \]


    \textbf{a) Alle Widerstände sind gleich groß: $R_1 = R_2 = R_3 = R_F = 10\,\mathrm{k}\Omega$}

    Eingangsspannungen:
    \[
        U_\mathrm{E1} = 1\,\mathrm{V}, \quad U_\mathrm{E2} = 2\,\mathrm{V}, \quad U_\mathrm{E3} = 3\,\mathrm{V}
    \]

    Berechnung:
    \[
        U_\mathrm{A} = -\left( 1\,\mathrm{V} \cdot \frac{10\,\mathrm{k}\Omega}{10\,\mathrm{k}\Omega} + 2\,\mathrm{V} \cdot \frac{10\,\mathrm{k}\Omega}{10\,\mathrm{k}\Omega} + 3\,\mathrm{V} \cdot \frac{10\,\mathrm{k}\Omega}{10\,\mathrm{k}\Omega} \right)
    \]
    \[
        U_A = -\left( 1\,\mathrm{V} + 2\,\mathrm{V} + 3\,\mathrm{V} \right) = -6\,\mathrm{V}
    \]


    \textbf{b) Geänderte Widerstände: $R_\mathrm{F} = 10\,\mathrm{k}\Omega$, $R_1 = 2\,\mathrm{k}\Omega$, $R_2 = 5\,\mathrm{k}\Omega$, $R_3 = 10\,\mathrm{k}\Omega$}

    Eingangsspannungen:
    \[
        U_\mathrm{E1} = 1\,\mathrm{V}, \quad U_\mathrm{E2} = 2\,\mathrm{V}, \quad U_\mathrm{E3} = 3\,\mathrm{V}
    \]

    Berechnung:
    \[
        U_\mathrm{A} = -\left( 1\,\mathrm{V} \cdot \frac{10\,\mathrm{k}\Omega}{2\,\mathrm{k}\Omega} + 2\,\mathrm{V} \cdot \frac{10\,\mathrm{k}\Omega}{5\,\mathrm{k}\Omega} + 3\,\mathrm{V} \cdot \frac{10\,\mathrm{k}\Omega}{10\,\mathrm{k}\Omega} \right)
    \]
    \[
        U_\mathrm{A} = -\left( 1\,\mathrm{V} \cdot 5 + 2\,\mathrm{V} \cdot 2 + 3\,\mathrm{V} \cdot 1 \right)
    \]
    \[
        U_\mathrm{A} = -(5\,\mathrm{V} + 4\,\mathrm{V} + 3\,\mathrm{V}) = -12\,\mathrm{V}
    \]


    \textbf{c) Geänderte Widerstände: $R_\mathrm{F} = 10\,\mathrm{k}\Omega$, $R_1 = 10\,\mathrm{k}\Omega$, $R_2 = 5\,\mathrm{k}\Omega$, $R_3 = 2\,\mathrm{k}\Omega$}

    Eingangsspannungen:
    \[
        U_\mathrm{E1} = 1\,\mathrm{V}, \quad U_\mathrm{E2} = 2\,\mathrm{V}, \quad U_\mathrm{E3} = 3\,\mathrm{V}
    \]

    Berechnung:
    \[
        U_\mathrm{A} = -\left( 1\,\mathrm{V} \cdot \frac{10\,\mathrm{k}\Omega}{10\,\mathrm{k}\Omega} + 2\,\mathrm{V} \cdot \frac{10\,\mathrm{k}\Omega}{5\,\mathrm{k}\Omega} + 3\,\mathrm{V} \cdot \frac{10\,\mathrm{k}\Omega}{2\,\mathrm{k}\Omega} \right)
    \]
    \[
        U_\mathrm{A} = -\left( 1\,\mathrm{V} \cdot 1 + 2\,\mathrm{V} \cdot 2 + 3\,\mathrm{V} \cdot 5 \right)
    \]
    \[
        U_\mathrm{A} = -(1\,\mathrm{V} + 4\,\mathrm{V} + 15\,\mathrm{V}) = -20\,\mathrm{V}
    \]

    Da die Betriebsspannung $\pm 15\,\mathrm{V}$ beträgt, wird der Ausgang auf $-15\,\mathrm{V}$ begrenzt.


    \textbf{Zusammenfassung:}
    \begin{itemize}
        \item a) $U_\mathrm{A} = -6\,\mathrm{V}$
        \item b) $U_\mathrm{A} = -12\,\mathrm{V}$
        \item c) $U_\mathrm{A} = -15\,\mathrm{V}$ (Begrenzung durch Betriebsspannung)
    \end{itemize}

    Siehe: Abschnitt \ref{fig: Verschaltung Verstaerker} und Tabelle \ref{tab:Grundschaltungen}
}